\section{Outcome (a): Sets, sample spaces, events, and the axioms of probability}

\subsection{What the outcome asks}

The exam does not ask you to recite definitions. It tests whether you can convert
a word problem into set notation, fill a two-event or three-event Venn diagram from
partial information, and apply the axioms and their direct consequences (complement
rule, inclusion-exclusion, monotonicity) to produce a number. The most common
question gives four to seven proportions of policyholders holding various coverages
and asks for the proportion in one region of the diagram or in a union of regions
(``none'', ``exactly one'', ``at most one'', ``auto but not home''). A second type
gives two probabilities with no other information and asks for the largest or
smallest value a third probability can take; that is solved with the axioms alone.
The arithmetic is short; points are lost in the translation from English to sets.

\subsection{Definitions}

\paragraph{Sample space, outcomes, events.}
The \emph{sample space} $S$ of a random experiment (select a policyholder at
random, observe next year's claims) is the set of all possible results; each element of $S$ is
an \emph{outcome}. An \emph{event} is a subset of $S$. $A$ \emph{occurs} when the observed outcome lies in $A$. $S$ is the certain
event and $\emptyset$ the impossible event.

\paragraph{Set operations.}
For events $A, B \subseteq S$:
\begin{itemize}[nosep]
  \item Union $A \cup B$: outcomes in $A$ or $B$ or both. English: ``$A$ or $B$'',
        ``at least one of $A$, $B$''.
  \item Intersection $A \cap B$ (also written $AB$): outcomes in both. English:
        ``$A$ and $B$'', ``both''.
  \item Complement $\cc{A} = S \setminus A$: outcomes not in $A$. English: ``not $A$''.
  \item Difference $A \setminus B = A \cap \cc{B}$: outcomes in $A$ but not in $B$.
        English: ``$A$ but not $B$'', ``$A$ only'' (when only two events are in play).
  \item $A$ and $B$ are \emph{disjoint} (mutually exclusive) when $A \cap B = \emptyset$;
        $A \subseteq B$ when every outcome of $A$ is in $B$, so $A$ occurring forces $B$.
\end{itemize}
Union and intersection distribute over each other:
$A \cap (B \cup C) = (A\cap B) \cup (A \cap C)$. Any event $A$ splits as the
disjoint union $A = (A \cap B) \cup (A \cap \cc{B})$; this identity is used constantly.

\begin{keyfact}[De Morgan's laws]
\[
  \cc{(A \cup B)} = \cc{A} \cap \cc{B},
  \qquad
  \cc{(A \cap B)} = \cc{A} \cup \cc{B}.
\]
In words: ``neither $A$ nor $B$'' is the complement of ``at least one'';
``not both'' is the complement of ``both''. The same pattern holds for any number
of events: the complement of a union is the intersection of the complements, and
the complement of an intersection is the union of the complements.
\end{keyfact}

\paragraph{Set functions and the collection of events.}
A \emph{set function} is a function whose inputs are sets. The number of elements
$|A|$, the length of an interval, and probability are all set functions.
The domain of a probability is a collection $\mathcal{F}$ of subsets of $S$, the
\emph{events}. $\mathcal{F}$ must contain $S$, must contain $\cc{A}$ whenever it
contains $A$, and must contain $A_1 \cup A_2 \cup \cdots$ whenever it contains
each $A_i$; a collection with these closure properties is a \emph{sigma-field}
(or sigma-algebra). When $S$ is finite, $\mathcal{F}$ is every subset of $S$; on the exam, every set
you can describe is an event.

\paragraph{Probability as a set function.}
A \emph{probability} is a set function $\Prob : \mathcal{F} \to \mathbb{R}$
satisfying three axioms (Kolmogorov, 1933).

\begin{keyfact}[The three axioms of probability]
\begin{enumerate}[nosep,label=\textbf{A\arabic*.}]
  \item Nonnegativity: $\Prob(A) \ge 0$ for every event $A$.
  \item Normalisation: $\Prob(S) = 1$.
  \item Countable additivity: if $A_1, A_2, \ldots$ are pairwise disjoint
        ($A_i \cap A_j = \emptyset$ for $i \ne j$), then
        \[
          \Prob\bigl(A_1 \cup A_2 \cup \cdots\bigr) = \Prob(A_1) + \Prob(A_2) + \cdots .
        \]
\end{enumerate}
\end{keyfact}

Everything else in this outcome is a consequence of the three axioms.

\begin{keyfact}[Consequences of the axioms]
\begin{enumerate}[nosep,label=\textbf{C\arabic*.}]
  \item $\Prob(\emptyset) = 0$.
  \item Finite additivity: for pairwise disjoint $A_1, \ldots, A_n$,
        $\Prob(A_1 \cup \cdots \cup A_n) = \sum_{i=1}^{n} \Prob(A_i)$.
  \item Complement rule: $\Prob(\cc{A}) = 1 - \Prob(A)$.
  \item Monotonicity: if $A \subseteq B$ then $\Prob(A) \le \Prob(B)$ and
        $\Prob(B \setminus A) = \Prob(B) - \Prob(A)$.
  \item $0 \le \Prob(A) \le 1$.
  \item Difference rule: $\Prob(A \cap \cc{B}) = \Prob(A) - \Prob(A \cap B)$.
  \item Inclusion-exclusion, two events:
        \[
          \Prob(A \cup B) = \Prob(A) + \Prob(B) - \Prob(A \cap B).
        \]
  \item Inclusion-exclusion, three events:
        \begin{align*}
          \Prob(A \cup B \cup C) = {}& \Prob(A) + \Prob(B) + \Prob(C) \\
          & - \Prob(A\cap B) - \Prob(A \cap C) - \Prob(B \cap C) \\
          & + \Prob(A \cap B \cap C).
        \end{align*}
  \item Boole's inequality: $\Prob(A_1 \cup \cdots \cup A_n) \le \sum_{i=1}^{n}\Prob(A_i)$.
  \item Bonferroni's inequality: $\Prob(A \cap B) \ge \Prob(A) + \Prob(B) - 1$.
\end{enumerate}
\end{keyfact}

Each consequence comes from writing a set as a disjoint union and applying A3:
$S = A \cup \cc{A}$ gives C3; $B = A \cup (B \setminus A)$ gives C4 and C5;
$A = (A\cap B) \cup (A\cap\cc{B})$ gives C6; $A \cup B = A \cup (B \cap \cc{A})$
gives C7, and C7 applied twice gives C8. When the exam asks which statements
``must be true'', the answer is one of C1--C10 or an axiom.

\paragraph{Identities from adding two inclusion-exclusion equations.}
Some questions give probabilities of unions that involve complements, such as
$\Prob(A \cup B)$ and $\Prob(A \cup \cc{B})$, and ask for $\Prob(A)$. Write
inclusion-exclusion for each union and add. Since $\Prob(B) + \Prob(\cc{B}) = 1$ and
$\Prob(A\cap B) + \Prob(A \cap \cc{B}) = \Prob(A)$ (C6), the sum collapses.

\begin{keyfact}[Union-with-complement identities]
\begin{align*}
  \Prob(A \cup B) + \Prob(A \cup \cc{B}) &= 1 + \Prob(A), \\
  \Prob(A \cup B) + \Prob(\cc{A} \cup B) &= 1 + \Prob(B), \\
  \Prob(A \cup \cc{B}) &= 1 - \Prob(\cc{A} \cap B), \\
  \Prob(\cc{A} \cup \cc{B}) &= 1 - \Prob(A \cap B), \\
  \Prob(A \cup B) + \Prob(\cc{A} \cup \cc{B}) &= 1 + \Prob(\text{exactly one of } A, B), \\
  \Prob(A \cup \cc{B}) + \Prob(\cc{A} \cup B) &= 2 - \Prob(\text{exactly one of } A, B).
\end{align*}
\end{keyfact}

The third and fourth lines are De Morgan plus the complement rule: $A \cup \cc{B}$
is the complement of $\cc{A} \cap B$, so its probability is one minus the
``$B$ but not $A$'' region. Every identity in the box is also a statement about the
$2\times 2$ table: $A \cup \cc{B}$ is three of the four cells (all but
$\cc{A}\cap B$), so $\Prob(A \cup B) + \Prob(A \cup \cc{B})$ counts the two $A$ cells
twice and the two $\cc{A}$ cells once, which is $\Prob(A) + 1$. When an identity is not
remembered, fill the table with unknowns and add cells.

\subsection{Venn diagram method}

A Venn diagram partitions $S$ into disjoint regions, so the probability of any
event is the sum of the regions it contains (C2). Name the events, fill the regions
from the innermost intersection outward, then read off the requested set as a sum
of regions.

\paragraph{Two events (four regions).}
The four regions are $A\cap B$, $A \cap \cc{B}$, $\cc{A} \cap B$, and
$\cc{A} \cap \cc{B}$. A $2\times 2$ table is the fastest bookkeeping; row and column
totals are the marginal probabilities, and the grand total is $1$.
\begin{center}
\begin{tabular}{@{}lccc@{}}
\toprule
 & $B$ & $\cc{B}$ & Total \\
\midrule
$A$      & $\Prob(A \cap B)$        & $\Prob(A \cap \cc{B})$        & $\Prob(A)$ \\
$\cc{A}$ & $\Prob(\cc{A} \cap B)$   & $\Prob(\cc{A} \cap \cc{B})$   & $\Prob(\cc{A})$ \\
\midrule
Total    & $\Prob(B)$               & $\Prob(\cc{B})$               & $1$ \\
\bottomrule
\end{tabular}
\end{center}
Any three independent pieces of information determine the whole table.

\paragraph{Three events (eight regions).}
Fill in the order shown. Each region's probability is obtained by subtracting
regions already known from a given total.
\begin{center}
\begin{tabular}{@{}clll@{}}
\toprule
Order & Region & Description & How to compute \\
\midrule
1 & $A \cap B \cap C$                   & all three          & given, or found last from the total \\
2 & $A \cap B \cap \cc{C}$              & $A$ and $B$ only   & $\Prob(A\cap B) - \Prob(A\cap B\cap C)$ \\
3 & $A \cap \cc{B} \cap C$              & $A$ and $C$ only   & $\Prob(A\cap C) - \Prob(A\cap B\cap C)$ \\
4 & $\cc{A} \cap B \cap C$              & $B$ and $C$ only   & $\Prob(B\cap C) - \Prob(A\cap B\cap C)$ \\
5 & $A \cap \cc{B} \cap \cc{C}$         & $A$ only           & $\Prob(A) - (\text{regions } 1,2,3)$ \\
6 & $\cc{A} \cap B \cap \cc{C}$         & $B$ only           & $\Prob(B) - (\text{regions } 1,2,4)$ \\
7 & $\cc{A} \cap \cc{B} \cap C$         & $C$ only           & $\Prob(C) - (\text{regions } 1,3,4)$ \\
8 & $\cc{A} \cap \cc{B} \cap \cc{C}$    & none               & $1 - (\text{regions } 1\text{--}7)$ \\
\bottomrule
\end{tabular}
\end{center}
When the problem gives ``none'' instead of ``all three'', the same table is solved
in reverse: the sum of regions 1--7 equals $1 - \Prob(\text{none})$, which with
inclusion-exclusion determines $\Prob(A \cap B \cap C)$.

\paragraph{Binary classifications.}
A problem that classifies each policyholder as young or old, male or female, and
married or single is a three-event problem in disguise: pick one label from each
pair as the event ($Y$, $M$, $R$ for young, male, married) and the other label is its
complement. ``Young, female, single'' is the region $Y \cap \cc{M} \cap \cc{R}$,
which is ``$Y$ only'' in the eight-region table. Counts work exactly like
probabilities; the grand total replaces $1$.

\paragraph{A disjoint pair inside a three-event problem.}
When the problem states that two of the three events are mutually exclusive, say
$B \cap C = \emptyset$, two regions vanish: ``$B$ and $C$ only'' and ``all three''.
Six regions remain: $A$ only, $B$ only, $C$ only, $A \cap B$, $A \cap C$, and none.
Name them with letters, translate every sentence of the problem into a linear
equation in those letters, and solve. Useful translations:
\begin{itemize}[nosep]
  \item ``exactly two of the three'' is the sum $(A\cap B) + (A \cap C)$;
  \item ``exactly one'' is $(A \text{ only}) + (B \text{ only}) + (C \text{ only})$;
  \item ``$B$ but not $A$'' is $B$ only (because $B \cap C$ is empty);
  \item ``not $A$'' is $(B \text{ only}) + (C \text{ only}) + (\text{none})$;
  \item ``twice as many have $B$ as have $C$'' is $(B\text{ only}) + (A \cap B) = 2\bigl[(C \text{ only}) + (A \cap C)\bigr]$.
\end{itemize}
The total of the six regions is $1$ (or the population size). Five independent
facts plus the total determine all six regions; the exam usually gives exactly
enough, and the order of solving is: everything outside $A$ first, then the
``exactly two'' total, then the ratio.

\paragraph{A fourth set disjoint from the other three.}
A survey may add a fourth category $D$ with the statement that no member of $D$
belongs to $A$, $B$, or $C$. Then $D$ is disjoint from $A \cup B \cup C$, and
by finite additivity
\[
  \Prob(A \cup B \cup C \cup D) = \Prob(A \cup B \cup C) + \Prob(D),
\]
so ``none of the four'' is $1 - \Prob(A\cup B\cup C) - \Prob(D)$. Run the
three-event inclusion-exclusion for $A$, $B$, $C$, then add $D$ as a ninth region.
Since $D$ overlaps nothing, ``exactly one of the four'' is ``exactly one of the
three'' plus $\Prob(D)$, and ``exactly two of the four'' equals ``exactly two of the
three''.

\paragraph{From English to sets.}
The two tables below are the translation dictionary; every phrase in them appears
in exam questions. Write $p_A = \Prob(A)$, $p_B = \Prob(B)$, $p_{AB} = \Prob(A \cap B)$.
\begin{center}
\begin{tabular}{@{}lll@{}}
\toprule
English phrase (two events) & Set expression & Probability \\
\midrule
$A$ or $B$; at least one            & $A \cup B$                                 & $p_A + p_B - p_{AB}$ \\
$A$ and $B$; both                   & $A \cap B$                                 & $p_{AB}$ \\
not $A$                             & $\cc{A}$                                   & $1 - p_A$ \\
neither $A$ nor $B$                 & $\cc{A} \cap \cc{B} = \cc{(A\cup B)}$      & $1 - p_A - p_B + p_{AB}$ \\
not both; at most one               & $\cc{A} \cup \cc{B} = \cc{(A\cap B)}$      & $1 - p_{AB}$ \\
$A$ but not $B$; $A$ only           & $A \cap \cc{B}$                            & $p_A - p_{AB}$ \\
exactly one of $A$, $B$             & $(A\cap\cc{B}) \cup (\cc{A}\cap B)$        & $p_A + p_B - 2p_{AB}$ \\
\bottomrule
\end{tabular}
\end{center}
For three events write $\Sigma_1 = \Prob(A)+\Prob(B)+\Prob(C)$,
$\Sigma_2 = \Prob(A\cap B)+\Prob(A\cap C)+\Prob(B\cap C)$, and
$T = \Prob(A\cap B\cap C)$. Region numbers refer to the eight-region table.
\begin{center}
\begin{tabular}{@{}lll@{}}
\toprule
English phrase (three events) & Set expression & Probability \\
\midrule
at least one                        & $A \cup B \cup C$                          & $\Sigma_1 - \Sigma_2 + T$ \\
none                                & $\cc{A}\cap\cc{B}\cap\cc{C}$               & $1 - \Sigma_1 + \Sigma_2 - T$ \\
all three                           & $A \cap B \cap C$                          & $T$ \\
$A$ only                            & $A \cap \cc{B} \cap \cc{C}$                & $\Prob(A) - \Prob(AB) - \Prob(AC) + T$ \\
exactly one                         & regions 5, 6, 7                            & $\Sigma_1 - 2\Sigma_2 + 3T$ \\
exactly two                         & regions 2, 3, 4                            & $\Sigma_2 - 3T$ \\
at least two                        & regions 1, 2, 3, 4                         & $\Sigma_2 - 2T$ \\
at most one                         & complement of ``at least two''             & $1 - \Sigma_2 + 2T$ \\
\bottomrule
\end{tabular}
\end{center}
Every three-event formula is a sum of regions from the eight-region table; when in
doubt, fill the regions and add.

\subsection{Countable additivity and geometric series}

Axiom A3 covers infinite sequences of disjoint events, and the exam tests it with
regions whose areas or probabilities form a geometric progression: the $i$th region
has probability $r^i$ (or $k r^i$) for $i = 1, 2, \ldots$, with $0 < r < 1$. The regions
are disjoint, so the probability of hitting one of the first $N$ is the sum of the
first $N$ terms, and the probability of hitting any of infinitely many is the sum of
the series.

\begin{keyfact}[Geometric sums for disjoint regions]
For $0 < r < 1$,
\[
  \sum_{i=1}^{N} r^i = \frac{r\,(1 - r^N)}{1 - r},
  \qquad
  \sum_{i=1}^{\infty} r^i = \frac{r}{1 - r},
  \qquad
  \sum_{i=0}^{\infty} r^i = \frac{1}{1 - r}.
\]
The infinite sum must not exceed $1$ (A3 with A2), so a design with areas $r^i$,
$i \ge 1$, is admissible only when $r \le 1/2$. If the union of the regions is the
whole space, the constant $k$ in $\Prob(A_i) = k r^i$ is fixed by
$k\,r/(1-r) = 1$.
\end{keyfact}

The probability of missing all $N$ regions is $1 - r(1 - r^N)/(1-r)$, which decreases
to the floor $1 - r/(1-r) = (1 - 2r)/(1-r)$ as $N \to \infty$. A ``minimum $N$ such
that the miss probability is below $p$'' question is solved by setting the closed
form below $p$, isolating $r^N$, and taking logarithms:
\[
  r^N < c \iff N > \frac{\ln c}{\ln r}
  \qquad (\text{the inequality flips because } \ln r < 0),
\]
then rounding up to the next integer. If $p$ is at or below the floor, no finite
$N$ works, and no infinite design works either.

\subsection{Testing whether two described events are equal}

A question may define $F$ and $G$ in words, list five events in words, and ask how
many of the five equal $F \cap G$ (or $F \cup G$). Reading alone is unreliable; the
descriptions are chosen to differ by one boundary case. The systematic method:
\begin{enumerate}[nosep]
  \item Choose coordinates for an outcome. For claims in two years, an outcome is
        the pair $(N_1, N_2)$ of claim counts.
  \item Coarsen each coordinate to the fewest classes that every description in the
        list can still distinguish. If the descriptions mention ``none'', ``exactly
        one'', and ``two or more'', use the classes $0$, $1$, $2^{+}$.
  \item Draw the grid of classes and mark the cells that belong to the target event.
  \item For each listed description, mark its cells and compare. Equal events have
        identical cell sets; one differing cell settles the case.
\end{enumerate}
The listed descriptions are usually algebraic rewrites of the target, and each
rewrite is right or wrong by one of the set laws:
\begin{itemize}[nosep]
  \item De Morgan: ``not ($\cc{F}$ and $\cc{G}$)'' equals $F \cup G$; ``not
        ($\cc{F}$ or $\cc{G}$)'' equals $F \cap G$.
  \item Disjoint decomposition: $F \cup G = F \cup (\cc{F} \cap G)$.
  \item Absorption: $F \cup (F \cap H) = F$ and $F \cap (F \cup H) = F$; adding a
        subset of $F$ to a union with $F$ changes nothing.
  \item Distribution: $F \cap (G \cup H) = (F\cap G) \cup (F \cap H)$.
  \item Given one coordinate, a condition on a total becomes a condition on the
        other coordinate: with $N_1 = 1$ fixed, ``$N_1 + N_2 \ge 2$'' is
        ``$N_2 \ge 1$''.
\end{itemize}
A description that is a proper subset of the target (it adds a condition) or a
proper superset (it drops one) is not equal to it, even when it is ``almost'' the
same. Watch for ``more than one'' (at least two) versus ``one or more'' (at least one),
and for a condition on the two-year total that silently allows a year with zero.

\subsection{Worked examples}

\begin{example}[Three-event coverage survey, none of the three]
An insurance company surveys its policyholders about three types of coverage:
auto, homeowners, and life. The survey finds that 40\% have auto coverage, 35\%
have homeowners coverage, and 25\% have life coverage. It also finds that 15\%
have both auto and homeowners, 10\% have both auto and life, 8\% have both
homeowners and life, and 5\% have all three coverages.

Calculate the percentage of policyholders with none of the three coverages.

(A) 22\% \quad (B) 28\% \quad (C) 33\% \quad (D) 38\% \quad (E) 72\%
\end{example}
\begin{solution}
Let $A$, $H$, $L$ be the events that a randomly chosen policyholder has auto,
homeowners, and life coverage. ``None of the three'' is $\cc{(A\cup H\cup L)}$, so
by the complement rule and three-event inclusion-exclusion,
\[
  \Prob(A \cup H \cup L) = 0.40 + 0.35 + 0.25 - 0.15 - 0.10 - 0.08 + 0.05 = 0.72,
\]
and $\Prob(\text{none}) = 1 - 0.72 = 0.28$.

Answer: (B).
\end{solution}

\begin{example}[Exactly one of three]
An insurance agency's records show that, for a randomly selected customer,
\begin{itemize}[nosep]
  \item the probability of having auto coverage is $0.50$, homeowners coverage
        $0.40$, and life coverage $0.30$;
  \item the probability of having both auto and homeowners is $0.20$, both auto
        and life is $0.15$, and both homeowners and life is $0.10$;
  \item the probability of having all three is $0.05$.
\end{itemize}
Calculate the probability that a randomly selected customer has exactly one of the
three coverages.

(A) 0.35 \quad (B) 0.40 \quad (C) 0.45 \quad (D) 0.50 \quad (E) 0.55
\end{example}
\begin{solution}
Fill the eight regions in the standard order. All three: $0.05$.
Pairs only: auto--homeowners $0.20 - 0.05 = 0.15$; auto--life $0.15 - 0.05 = 0.10$;
homeowners--life $0.10 - 0.05 = 0.05$.
Singles only: auto $0.50 - 0.15 - 0.10 - 0.05 = 0.20$;
homeowners $0.40 - 0.15 - 0.05 - 0.05 = 0.15$; life $0.30 - 0.10 - 0.05 - 0.05 = 0.10$.
None: $1 - 0.80 = 0.20$.

Exactly one is the sum of the three ``only'' regions: $0.20 + 0.15 + 0.10 = 0.45$.

Formula check: $\Sigma_1 - 2\Sigma_2 + 3T = 1.20 - 2(0.45) + 3(0.05) = 1.20 - 0.90 + 0.15 = 0.45$.

Answer: (C).
\end{solution}

\begin{example}[Bounding a probability from the axioms alone]
For a randomly selected claim from an auto insurer's files, the probability that
the claim involves bodily injury is $0.60$ and the probability that it involves
property damage is $0.70$. No other information about the claims is available.

Determine the smallest possible value of the probability that a claim involves
both bodily injury and property damage.

(A) 0 \quad (B) 0.10 \quad (C) 0.30 \quad (D) 0.42 \quad (E) 0.60
\end{example}
\begin{solution}
Let $B$ be bodily injury and $D$ be property damage. Inclusion-exclusion gives
$\Prob(B \cap D) = \Prob(B) + \Prob(D) - \Prob(B \cup D) = 1.30 - \Prob(B \cup D)$.
Since $\Prob(B \cup D) \le 1$ (C5), $\Prob(B \cap D) \ge 1.30 - 1 = 0.30$. The
value $0.30$ is attained when $B \cup D = S$, so it is the minimum. (This is
Bonferroni's inequality, C10.)

The largest possible value is $\min(0.60, 0.70) = 0.60$, attained when
$B \subseteq D$ (C4). The choice $0.42 = (0.60)(0.70)$ assumes independence, which
is not given.

Answer: (C).
\end{solution}

\begin{example}[Complements and De Morgan]
A property insurer's records show that, for a randomly selected policyholder, the
probability of not having auto coverage is $0.65$, the probability of not having
homeowners coverage is $0.55$, and the probability of having neither coverage is
$0.35$.

Calculate the probability that a randomly selected policyholder has both auto and
homeowners coverage.

(A) 0.10 \quad (B) 0.15 \quad (C) 0.20 \quad (D) 0.35 \quad (E) 0.45
\end{example}
\begin{solution}
Let $A$ and $H$ be auto and homeowners coverage. Translate each statement.
$\Prob(\cc{A}) = 0.65$ so $\Prob(A) = 0.35$.
$\Prob(\cc{H}) = 0.55$ so $\Prob(H) = 0.45$.
``Neither'' is $\cc{A} \cap \cc{H} = \cc{(A \cup H)}$ by De Morgan, so
$\Prob(A \cup H) = 1 - 0.35 = 0.65$.
Then
\[
  \Prob(A \cap H) = \Prob(A) + \Prob(H) - \Prob(A \cup H) = 0.35 + 0.45 - 0.65 = 0.15 .
\]
Table check: $A \cap \cc{H} = 0.35 - 0.15 = 0.20$, $\cc{A} \cap H = 0.45 - 0.15 = 0.30$,
neither $= 0.35$; the four regions sum to $0.15 + 0.20 + 0.30 + 0.35 = 1.00$.

Answer: (B).
\end{solution}

\begin{example}[Translating ``only'', ``none'', ``no customer has both'', ``at most one'']
An insurance agency has 200 customers. Each customer has some combination of auto,
homeowners, and renters coverage, or none of them. The agency's records show:
\begin{itemize}[nosep]
  \item 120 customers have auto coverage;
  \item 60 customers have homeowners coverage and 50 have renters coverage; no
        customer has both homeowners and renters coverage;
  \item 65 customers have auto coverage only;
  \item 25 customers have none of the three coverages.
\end{itemize}
Calculate the number of customers who have at most one of the three coverages.

(A) 55 \quad (B) 65 \quad (C) 120 \quad (D) 145 \quad (E) 175
\end{example}
\begin{solution}
Let $A$, $H$, $R$ be the three coverages and work with counts out of 200.
``No customer has both homeowners and renters'' means $H \cap R = \emptyset$, so
$A \cap H \cap R$ is also empty and every two-coverage customer has auto plus one
other. ``Auto only'' is 65, so $120 - 65 = 55$ customers have auto plus another
coverage; these are exactly the two-coverage customers. Since $H$ and $R$ are
disjoint, $|H \cup R| = 60 + 50 = 110$, of whom 55 also have auto, leaving
$110 - 55 = 55$ with homeowners only or renters only. Exactly one coverage:
$65 + 55 = 120$. Check: $120 + 55 + 25 = 200$.

``At most one'' is none or exactly one: $25 + 120 = 145$.

Answer: (D).
\end{solution}

\begin{example}[Unions with complements]
For a randomly selected policy in a commercial lines portfolio, let $A$ be the event
that the policy carries general liability coverage and $B$ the event that it carries
property coverage. You are given
\[
  \Prob(A \cup B) = 0.78, \qquad \Prob(A \cup \cc{B}) = 0.67, \qquad \Prob(B) = 0.55 .
\]
Calculate the probability that the policy carries both coverages.

(A) 0.22 \qquad (B) 0.33 \qquad (C) 0.45 \qquad (D) 0.55 \qquad (E) 0.67
\end{example}
\begin{solution}
Write inclusion-exclusion for each union and add:
\begin{align*}
  \Prob(A \cup B) &= \Prob(A) + \Prob(B) - \Prob(A \cap B), \\
  \Prob(A \cup \cc{B}) &= \Prob(A) + \Prob(\cc{B}) - \Prob(A \cap \cc{B}), \\
  0.78 + 0.67 &= 2\Prob(A) + 1 - \bigl[\Prob(A\cap B) + \Prob(A \cap \cc{B})\bigr]
               = 2\Prob(A) + 1 - \Prob(A),
\end{align*}
so $\Prob(A) = 1.45 - 1 = 0.45$. Then
$\Prob(A \cap B) = \Prob(A) + \Prob(B) - \Prob(A \cup B) = 0.45 + 0.55 - 0.78 = 0.22$.

Shorter route: $A \cup \cc{B}$ is the complement of $\cc{A} \cap B$ (De Morgan), so
$\Prob(\cc{A} \cap B) = 1 - 0.67 = 0.33$, and
$\Prob(A \cap B) = \Prob(B) - \Prob(\cc{A}\cap B) = 0.55 - 0.33 = 0.22$.

Table check: $A\cap B = 0.22$, $A \cap \cc{B} = 0.45 - 0.22 = 0.23$,
$\cc{A} \cap B = 0.33$, neither $= 1 - 0.78 = 0.22$; the sum is $1.00$.

Answer: (A).
\end{solution}

\begin{example}[Disjoint pair inside a three-product profile]
An agent sells disability income insurance, term life insurance, and whole life
insurance. No client holds both term life and whole life. The agent's client
profile is:
\begin{enumerate}[nosep,label=(\roman*)]
  \item 20\% of clients hold none of the three products;
  \item 63\% of clients hold disability income insurance;
  \item three times as many clients hold term life as hold whole life;
  \item 23\% of clients hold exactly two of the three products;
  \item 12\% of clients hold term life but not disability income insurance.
\end{enumerate}
Calculate the percentage of clients who hold both disability income and whole life
insurance.

(A) 5\% \qquad (B) 8\% \qquad (C) 12\% \qquad (D) 15\% \qquad (E) 18\%
\end{example}
\begin{solution}
Let $D$, $T$, $W$ be the three products. Since $T \cap W = \emptyset$, the regions
$T \cap W$ only and $D \cap T \cap W$ are empty, leaving six regions: $d$ ($D$ only),
$t$ ($T$ only), $w$ ($W$ only), $dt$, $dw$, and $n$ (none). Translate:
\[
  n = 0.20, \quad d + dt + dw = 0.63, \quad t + dt = 3(w + dw), \quad dt + dw = 0.23,
  \quad t = 0.12 .
\]
Fact (v) gives $t$ directly because $T$ but not $D$ cannot include any $W$.
Outside $D$ lie $t$, $w$, and $n$, with total $1 - 0.63 = 0.37$, so
$w = 0.37 - 0.20 - 0.12 = 0.05$. Substitute into the ratio:
$0.12 + dt = 3(0.05 + dw) = 0.15 + 3\,dw$, so $dt = 0.03 + 3\,dw$. With
$dt + dw = 0.23$: $0.03 + 4\,dw = 0.23$, so $dw = 0.05$ and $dt = 0.18$.

Check: $d = 0.63 - 0.18 - 0.05 = 0.40$, and
$0.40 + 0.12 + 0.05 + 0.18 + 0.05 + 0.20 = 1.00$. Term life totals
$0.12 + 0.18 = 0.30$, whole life $0.05 + 0.05 = 0.10$, ratio $3$.

Answer: (A).
\end{solution}

\begin{example}[Three overlapping benefits and a fourth disjoint one]
A company offers its 250 employees four voluntary benefits: dental, vision, legal,
and a commuter plan. Enrollment records show:
\begin{enumerate}[nosep,label=(\roman*)]
  \item 96 are enrolled in dental, 58 in vision, and 40 in legal;
  \item 27 are enrolled in dental and vision, 18 in dental and legal, and 12 in
        vision and legal;
  \item 7 are enrolled in all three of dental, vision, and legal;
  \item 35 are enrolled in the commuter plan, and none of these is enrolled in
        dental, vision, or legal.
\end{enumerate}
Calculate the number of employees enrolled in exactly one of the four benefits.

(A) 71 \qquad (B) 101 \qquad (C) 114 \qquad (D) 136 \qquad (E) 179
\end{example}
\begin{solution}
Let $D$, $V$, $L$, $K$ be the four benefits. For the three overlapping ones,
$\Sigma_1 = 96 + 58 + 40 = 194$, $\Sigma_2 = 27 + 18 + 12 = 57$, $T = 7$.
Exactly one of $D$, $V$, $L$ is $\Sigma_1 - 2\Sigma_2 + 3T = 194 - 114 + 21 = 101$.
Region check: $D$ only $= 96 - 20 - 11 - 7 = 58$, $V$ only $= 58 - 20 - 5 - 7 = 26$,
$L$ only $= 40 - 11 - 5 - 7 = 17$, total $101$.

$K$ is disjoint from $D \cup V \cup L$, so every one of its 35 members has exactly one
benefit, and none of them changes the count for $D$, $V$, $L$. Exactly one of the
four: $101 + 35 = 136$.

For reference, $|D \cup V \cup L| = 194 - 57 + 7 = 144$, $|D\cup V\cup L\cup K| = 144 + 35 = 179$
by finite additivity, and none of the four $= 250 - 179 = 71$.

Answer: (D).
\end{solution}

\begin{example}[Countable additivity with geometric areas]
A carnival game uses a board of area 1 painted in red and blue regions. A
blindfolded player marks a point on the board and wins if the point lies in a blue
region. The board can be painted with any number of red regions, and the red regions
are designed so that the $i$th red region has area $(2/5)^i$, $i = 1, 2, \ldots$.

Calculate the minimum number of red regions that makes the probability of winning
less than 34\%.

(A) 2 \qquad (B) 3 \qquad (C) 4 \qquad (D) 5 \qquad (E) 6
\end{example}
\begin{solution}
With $N$ red regions, the red regions are disjoint, so by additivity the losing
probability is the sum of their areas and the winning probability is
\[
  w_N = 1 - \sum_{i=1}^{N} \Bigl(\frac{2}{5}\Bigr)^i
      = 1 - \frac{(2/5)\bigl(1 - (2/5)^N\bigr)}{1 - 2/5}
      = 1 - \frac{2}{3}\Bigl(1 - \Bigl(\frac{2}{5}\Bigr)^N\Bigr)
      = \frac{1}{3} + \frac{2}{3}\Bigl(\frac{2}{5}\Bigr)^N .
\]
Require $w_N < 0.34$:
\[
  \frac{2}{3}\Bigl(\frac{2}{5}\Bigr)^N < 0.34 - \frac{1}{3} = \frac{1}{150}
  \iff \Bigl(\frac{2}{5}\Bigr)^N < \frac{1}{100}
  \iff N > \frac{\ln(0.01)}{\ln(0.4)} = 5.03 .
\]
The smallest integer is $N = 6$. Check: $w_5 = 1/3 + (2/3)(0.01024) = 0.3402$, not
below $0.34$; $w_6 = 1/3 + (2/3)(0.004096) = 0.3361$.

The winning probability never drops below the floor $1/3$: the infinite series of red
areas sums to $(2/5)/(3/5) = 2/3$. A target of ``less than 33\%'' would have no
answer, with any finite or infinite number of red regions.

Answer: (E).
\end{solution}

\begin{example}[Which described events equal $F \cup G$]
A policyholder holds an auto policy for two years. Let $F$ be the event that the
policyholder has at most one accident in year one, and $G$ the event that the
policyholder has no accidents in year two. Consider the following events:
\begin{enumerate}[nosep,label=(\roman*)]
  \item The policyholder has at most one accident in the two-year period.
  \item It is not the case that the policyholder has two or more accidents in year
        one and at least one accident in year two.
  \item The policyholder has at most one accident in year one, or has two or more
        accidents in year one and no accidents in year two.
  \item The policyholder has at most one accident in year one, or has no accidents
        in either year.
  \item The policyholder has at most one accident in year one, or has more
        accidents in year one than in year two and no accidents in year two.
\end{enumerate}
Determine the number of events from the list that are the same as $F \cup G$.

(A) None \qquad (B) Exactly one \qquad (C) Exactly two \qquad (D) Exactly three \qquad (E) All
\end{example}
\begin{solution}
Let $N_1$ and $N_2$ be the accident counts in the two years. Every description
distinguishes only the classes $0$, $1$, and $2^{+}$ for each year, so the outcome
grid has nine cells. $F$ is the rows $N_1 = 0$ and $N_1 = 1$; $G$ is the column
$N_2 = 0$; $F \cup G$ is those two rows plus the cell $(2^{+}, 0)$, seven cells in all.
A bullet marks membership.
\begin{center}
\begin{tabular}{@{}lcccccc@{}}
\toprule
$(N_1, N_2)$ & $F \cup G$ & (i) & (ii) & (iii) & (iv) & (v) \\
\midrule
$(0,0)$          & $\bullet$ & $\bullet$ & $\bullet$ & $\bullet$ & $\bullet$ & $\bullet$ \\
$(0,1)$          & $\bullet$ & $\bullet$ & $\bullet$ & $\bullet$ & $\bullet$ & $\bullet$ \\
$(0,2^{+})$      & $\bullet$ &           & $\bullet$ & $\bullet$ & $\bullet$ & $\bullet$ \\
$(1,0)$          & $\bullet$ & $\bullet$ & $\bullet$ & $\bullet$ & $\bullet$ & $\bullet$ \\
$(1,1)$          & $\bullet$ &           & $\bullet$ & $\bullet$ & $\bullet$ & $\bullet$ \\
$(1,2^{+})$      & $\bullet$ &           & $\bullet$ & $\bullet$ & $\bullet$ & $\bullet$ \\
$(2^{+},0)$      & $\bullet$ &           & $\bullet$ & $\bullet$ &           & $\bullet$ \\
$(2^{+},1)$      &           &           &           &           &           &           \\
$(2^{+},2^{+})$  &           &           &           &           &           &           \\
\bottomrule
\end{tabular}
\end{center}
(i) is the event $N_1 + N_2 \le 1$, three cells; it omits $(0, 2^{+})$, which is in
$F$. Not equal.
(ii) is the complement of $\cc{F} \cap \cc{G}$, which by De Morgan is $F \cup G$. Equal.
(iii) is $F \cup (\cc{F} \cap G)$, the disjoint decomposition of $F \cup G$. Equal.
(iv) is $F \cup \{(0,0)\}$; the added cell is already in $F$, so this is just $F$,
which lacks $(2^{+}, 0)$. Not equal.
(v) is $F \cup \{N_1 > N_2,\ N_2 = 0\} = F \cup \{N_1 \ge 1,\ N_2 = 0\}$. The added
cells are $(1,0)$, already in $F$, and $(2^{+},0)$; the result is $F \cup G$. Equal.

Three of the five events equal $F \cup G$.

Answer: (D).
\end{solution}

\subsection{Traps}

\begin{trap}[``At least one'' versus ``exactly one'' versus ``at most one'']
``At least one of $A$, $B$'' is $A \cup B$ and includes the intersection. ``Exactly
one'' excludes the intersection: $\Prob(A) + \Prob(B) - 2\Prob(A\cap B)$. ``At most
one'' is the complement of ``both'': $1 - \Prob(A \cap B)$. For three events,
``at most one'' is not $1 - T$; it is $1 - (\Sigma_2 - 2T)$, the complement of ``at
least two''. Likewise ``neither $A$ nor $B$'' is $\cc{(A\cup B)}$ while ``not both''
is $\cc{(A \cap B)}$; De Morgan's laws tell you which is which. \end{trap}

\begin{trap}[$\Prob(A \cap \cc{B})$ is not $\Prob(A) - \Prob(B)$]
``$A$ but not $B$'' has probability $\Prob(A) - \Prob(A \cap B)$. Subtracting
$\Prob(B)$ instead removes outcomes that were never in $A$. The identity
$\Prob(B \setminus A) = \Prob(B) - \Prob(A)$ holds only when $A \subseteq B$.
\end{trap}

\begin{trap}[Assuming disjoint or independent when the problem does not say so]
If a question gives $\Prob(A)$ and $\Prob(B)$ and asks for $\Prob(A \cup B)$ with
no other information, the answer is not determined; it lies between
$\max(\Prob(A),\Prob(B))$ and $\min(1, \Prob(A)+\Prob(B))$. Adding the probabilities
assumes disjointness; multiplying them assumes independence. Use either only when
the problem states it or the structure forces it (a policyholder cannot be both
under 25 and over 65).
\end{trap}

\begin{trap}[``$A$ and $B$'' includes those who also have $C$]
In a three-event problem, ``15\% have both auto and homeowners'' is
$\Prob(A \cap H) = 0.15$, which includes the policyholders who also have life.
The region ``auto and homeowners only'' is $\Prob(A\cap H) - \Prob(A\cap H\cap L)$.
Conversely, when the problem says ``15\% have auto and homeowners but not life'',
it is giving the region directly. Confusing the two shifts the answer by a multiple of $T$.
\end{trap}

\begin{trap}[$\Prob(A \cup \cc{B})$ is not $1 - \Prob(A \cup B)$]
The complement of $A \cup B$ is $\cc{A} \cap \cc{B}$, not $A \cup \cc{B}$. The event
$A \cup \cc{B}$ is three of the four cells of the $2 \times 2$ table (everything except
$\cc{A} \cap B$), so its probability is $1 - \Prob(\cc{A} \cap B)$, and it is
typically large. When two unions involving complements are given, add the two
inclusion-exclusion equations rather than guessing a relation between them.
\end{trap}

\begin{trap}[An infinite union of disjoint regions has a floor, and a ceiling]
With region probabilities $r^i$, $i \ge 1$, the probability of missing every
region decreases toward $(1-2r)/(1-r)$, not toward $0$; a question that asks for a miss
probability below that floor has no answer. In the other direction, the series must
sum to at most $1$, so $r^i$ areas on a board of area $1$ require $r \le 1/2$. When
solving $r^N < c$, dividing by $\ln r$ (negative) reverses the inequality, and the
answer is the next integer above the bound, not the nearest integer.
\end{trap}

\begin{trap}[Coarsening the outcome grid too far]
The classes $0$, $1$, $2^{+}$ are enough when every description compares a count
to a fixed number. A description that compares the two years to each other (``more
accidents in year one than in year two'') can split the $2^{+}$ class: the outcomes
$(3,2)$ and $(2,3)$ both sit in the cell $(2^{+}, 2^{+})$ but fall on opposite sides of
the comparison. Either refine the classes or test the ambiguous cell with explicit
outcomes before declaring two events equal.
\end{trap}

\subsection{Practice problems}

\begin{practice}
A survey of homeowners insurance policyholders finds that 50\% have a security
system, 35\% have a fire sprinkler system, and 25\% have a backup generator.
Further, 15\% have both a security system and a sprinkler system, 10\% have both a
security system and a generator, 5\% have both a sprinkler system and a generator,
and 2\% have all three.

Calculate the percentage of policyholders with none of the three.

(A) 12\% \quad (B) 18\% \quad (C) 20\% \quad (D) 28\% \quad (E) 32\%
\end{practice}

\begin{practice}
For a randomly selected policyholder, let $A$ be the event that the policyholder
files a claim in the first policy year and $B$ the event that the policyholder
files a claim in the second policy year. You are given
$\Prob(\cc{A}) = 0.60$, $\Prob(B) = 0.35$, and $\Prob(A \cup B) = 0.60$.

Calculate the probability that the policyholder files a claim in exactly one of the
two years.

(A) 0.30 \quad (B) 0.45 \quad (C) 0.55 \quad (D) 0.60 \quad (E) 0.75
\end{practice}

\begin{practice}
For a randomly selected claim, the probability that the claim involves a rental
vehicle is $0.75$ and the probability that it involves a towing charge is $0.65$.
No other information is available.

Determine the largest possible value of the probability that the claim involves
neither a rental vehicle nor a towing charge.

(A) 0 \quad (B) 0.0875 \quad (C) 0.25 \quad (D) 0.35 \quad (E) 0.40
\end{practice}

\begin{practice}
An insurer classifies each policy by whether it has a young driver ($Y$), a
sports car ($S$), and an accident in the past three years ($C$). For a randomly
selected policy,
\begin{gather*}
  \Prob(Y) = 0.30,\quad \Prob(S) = 0.25,\quad \Prob(C) = 0.20,\\
  \Prob(Y\cap S) = 0.10,\quad \Prob(Y \cap C) = 0.08,\quad \Prob(S \cap C) = 0.06,\quad
  \Prob(Y \cap S \cap C) = 0.04 .
\end{gather*}
Calculate the probability that a randomly selected policy has at most one of the
three characteristics.

(A) 0.16 \quad (B) 0.24 \quad (C) 0.76 \quad (D) 0.84 \quad (E) 0.90
\end{practice}

\begin{practice}
The probability that a randomly selected auto policyholder lacks at least one of
collision coverage and comprehensive coverage is $0.90$. The probability of having
collision coverage is $0.45$ and the probability of having comprehensive coverage
is $0.30$.

Calculate the probability that the policyholder has at least one of the two
coverages.

(A) 0.10 \quad (B) 0.35 \quad (C) 0.55 \quad (D) 0.65 \quad (E) 0.75
\end{practice}

\begin{practice}
Of 500 auto policyholders, 280 have collision coverage, 340 have comprehensive
coverage, and 60 have neither.

Calculate the number of policyholders who have collision coverage only.

(A) 80 \quad (B) 100 \quad (C) 160 \quad (D) 180 \quad (E) 220
\end{practice}

\begin{practice}
For a randomly selected policyholder, let $A$ be the event that the policyholder
files a claim this year and $B$ the event that the policyholder renews at year end.
You are given
\[
  \Prob(A \cup B) = 0.85, \qquad \Prob(A \cup \cc{B}) = 0.75, \qquad \Prob(\cc{A} \cup B) = 0.90 .
\]
Calculate the probability that at most one of the events $A$ and $B$ occurs.

(A) 0.15 \qquad (B) 0.25 \qquad (C) 0.35 \qquad (D) 0.40 \qquad (E) 0.50
\end{practice}

\begin{practice}
An insurance agent sells auto, homeowners, and renters insurance. No client holds
both homeowners and renters insurance. The agent's client profile is:
\begin{enumerate}[nosep,label=(\roman*)]
  \item 12\% of clients hold none of the three products;
  \item 43\% of clients hold exactly one of the three products;
  \item 73\% of clients hold auto insurance;
  \item 5\% of clients hold renters insurance but not auto insurance;
  \item twice as many clients hold both auto and homeowners insurance as hold both
        auto and renters insurance.
\end{enumerate}
Calculate the percentage of clients who hold homeowners insurance.

(A) 20\% \qquad (B) 28\% \qquad (C) 30\% \qquad (D) 38\% \qquad (E) 40\%
\end{practice}

\begin{practice}
A bank surveys 400 customers about four products: a checking account, a savings
account, a brokerage account, and a certificate of deposit. The survey finds:
\begin{enumerate}[nosep,label=(\roman*)]
  \item 130 hold a checking account, 100 a savings account, and 70 a brokerage
        account;
  \item 40 hold checking and savings, 30 hold checking and brokerage, and 20 hold
        savings and brokerage;
  \item 60 hold a certificate of deposit, and none of these holds a checking,
        savings, or brokerage account;
  \item 115 hold none of the four products.
\end{enumerate}
Calculate the number of customers who hold all three of the checking, savings, and
brokerage accounts.

(A) 15 \qquad (B) 25 \qquad (C) 30 \qquad (D) 45 \qquad (E) 60
\end{practice}

\begin{practice}
Each incoming claim is assigned to exactly one adjuster from the sequence
$A_1, A_2, A_3, \ldots$, and the probability that a claim is assigned to adjuster
$A_i$ is $k\,(3/4)^i$ for $i = 1, 2, \ldots$, where $k$ is a constant.

Calculate the smallest $n$ such that the probability that a claim is assigned to one
of the first $n$ adjusters is at least $0.95$.

(A) 9 \qquad (B) 10 \qquad (C) 11 \qquad (D) 12 \qquad (E) 13
\end{practice}

\begin{practice}
A policyholder holds a policy for two years. Let $F$ be the event that the
policyholder has no claims in year one and $G$ the event that the policyholder has
at least two claims in year two. Consider the following events:
\begin{enumerate}[nosep,label=(\roman*)]
  \item The policyholder has at least one claim in year one and at most one claim in
        year two.
  \item It is not the case that the policyholder has no claims in year one or has at
        least two claims in year two.
  \item The policyholder has at least one claim in year one, and has either no claims
        or exactly one claim in year two.
  \item The policyholder has exactly one claim in year one and at most one claim in
        year two, or has two or more claims in year one and at most one claim in
        year two.
  \item The policyholder has at least one claim in year one and at most one claim in
        year two, or has exactly one claim in year one and no claims in year two.
\end{enumerate}
Determine the number of events from the list that are the same as $\cc{(F \cup G)}$.

(A) None \qquad (B) Exactly one \qquad (C) Exactly two \qquad (D) Exactly three \qquad (E) All
\end{practice}

\begin{practice}
An auto insurer has 5{,}000 policies. Each policy is classified as urban or rural,
as covering a sedan or a truck, and as commuting use or pleasure use. Of these
policies, 2{,}100 are urban, 2{,}800 cover sedans, and 3{,}200 are for commuting
use. Further, 1{,}150 are urban sedans, 1{,}900 are commuting sedans, and 1{,}350
are urban commuting policies. Finally, 600 policies are rural trucks for pleasure use.

Calculate the number of policies that are urban trucks for pleasure use.

(A) 300 \qquad (B) 450 \qquad (C) 650 \qquad (D) 700 \qquad (E) 950
\end{practice}

\subsection{Answers to practice problems}

\begin{enumerate}[nosep,label=\textbf{\arabic*.}]
  \item \textbf{(B)}. $\Prob(\text{at least one}) = 1.10 - 0.30 + 0.02 = 0.82$,
        so none is $0.18$.
  \item \textbf{(B)}. $\Prob(A) = 0.40$; $\Prob(A\cap B) = 0.40 + 0.35 - 0.60 = 0.15$;
        exactly one $= 0.40 + 0.35 - 2(0.15) = 0.45$ (equivalently $0.60 - 0.15$).
  \item \textbf{(C)}. Neither $= 1 - \Prob(R \cup T)$, and $\Prob(R \cup T) \ge \max(0.75, 0.65) = 0.75$
        by monotonicity, so neither $\le 0.25$, attained when $T \subseteq R$.
        (0.0875 assumes independence, which is not given.)
  \item \textbf{(D)}. At least two $= \Sigma_2 - 2T = 0.24 - 0.08 = 0.16$; at most one $= 0.84$.
  \item \textbf{(D)}. ``Lacks at least one'' is $\cc{(C \cap M)}$, so $\Prob(C \cap M) = 0.10$;
        $\Prob(C \cup M) = 0.45 + 0.30 - 0.10 = 0.65$.
  \item \textbf{(B)}. At least one $= 440$; both $= 280 + 340 - 440 = 180$;
        collision only $= 280 - 180 = 100$.
  \item \textbf{(E)}. $\Prob(A) = 0.85 + 0.75 - 1 = 0.60$; $\Prob(B) = 0.85 + 0.90 - 1 = 0.75$;
        $\Prob(A \cap B) = 0.60 + 0.75 - 0.85 = 0.50$; at most one $= 1 - 0.50 = 0.50$.
        (Neither $= 0.15$; exactly one $= 0.35$.)
  \item \textbf{(E)}. Exactly two $= 1 - 0.12 - 0.43 = 0.45$; auto only $= 0.73 - 0.45 = 0.28$;
        homeowners only $+$ renters only $= 0.43 - 0.28 = 0.15$, so homeowners only $= 0.10$;
        auto-renters $= 0.15$ and auto-homeowners $= 0.30$ from the ratio;
        homeowners $= 0.10 + 0.30 = 0.40$.
  \item \textbf{(A)}. The certificate holders are disjoint from the other three, so
        $|C \cup S \cup B| = 400 - 115 - 60 = 225 = 300 - 90 + T$, giving $T = 15$.
  \item \textbf{(C)}. The union of all $A_i$ is $S$, so $k(3/4)/(1/4) = 3k = 1$ and $k = 1/3$.
        $\Prob(A_1 \cup \cdots \cup A_n) = 1 - (3/4)^n \ge 0.95$ requires
        $n \ge \ln(0.05)/\ln(0.75) = 10.41$, so $n = 11$.
  \item \textbf{(E)}. The target is $\cc{F} \cap \cc{G} = \{N_1 \ge 1, N_2 \le 1\}$.
        (i) is the direct statement; (ii) is De Morgan; (iii) splits $N_2 \le 1$ into
        $N_2 = 0$ or $N_2 = 1$; (iv) splits $N_1 \ge 1$ into $N_1 = 1$ or $N_1 \ge 2$;
        (v) adds the cell $(1, 0)$, already in the target, so absorption leaves it unchanged.
  \item \textbf{(A)}. Union of urban, sedan, commuting $= 5000 - 600 = 4400 = 8100 - 4400 + T$,
        so $T = 700$; urban only $= 2100 - 1150 - 1350 + 700 = 300$.
\end{enumerate}

\subsection{Sample exam cross-reference}

After studying this section, test on the following questions from the sample exam.
\begin{itemize}[nosep]
  \item Bank 1, Q1: three-event inclusion-exclusion, none of the three (Example 1, Practice 1).
  \item Bank 1, Q3: add $\Prob(A \cup B)$ and $\Prob(A \cup \cc{B})$ to isolate $\Prob(A)$ (Example 6, Practice 7).
  \item Bank 1, Q5: three binary classifications as a three-event count; ``young, female, single'' is a single region (Example 2, Practice 12).
  \item Bank 1, Q87: disjoint pair inside three products, six-region linear system with an ``exactly two'' constraint and a ratio (Example 7, Practice 8).
  \item Bank 1, Q100: three overlapping sets plus a fourth disjoint from them; finite additivity for the union (Example 8, Practice 9).
  \item Bank 1, Q108: countable additivity over disjoint regions with geometric areas; minimum $N$ by logarithms (Example 9, Practice 10).
  \item Bank 1, Q153: event equality tested on an outcome grid; De Morgan, decomposition, absorption (Example 10, Practice 11).
\end{itemize}

\subsection{One-page summary}

\begin{itemize}[nosep]
  \item Sample space $S$: all outcomes. Event: a subset of $S$. Probability: a set
        function on the events with $\Prob(A) \ge 0$, $\Prob(S) = 1$, and countable
        additivity over pairwise disjoint events.
  \item Consequences: $\Prob(\emptyset) = 0$; $\Prob(\cc{A}) = 1 - \Prob(A)$;
        $A \subseteq B \Rightarrow \Prob(A) \le \Prob(B)$; $0 \le \Prob(A) \le 1$;
        $\Prob(A \cap \cc{B}) = \Prob(A) - \Prob(A \cap B)$.
  \item De Morgan: $\cc{(A \cup B)} = \cc{A} \cap \cc{B}$ (``neither'');
        $\cc{(A \cap B)} = \cc{A} \cup \cc{B}$ (``not both'').
  \item Inclusion-exclusion: $\Prob(A \cup B) = \Prob(A) + \Prob(B) - \Prob(A \cap B)$;
        $\Prob(A\cup B\cup C) = \Sigma_1 - \Sigma_2 + T$.
  \item Exactly one of two: $\Prob(A) + \Prob(B) - 2\Prob(A\cap B)$.
        Exactly one of three: $\Sigma_1 - 2\Sigma_2 + 3T$.
        Exactly two of three: $\Sigma_2 - 3T$. At least two of three: $\Sigma_2 - 2T$.
  \item Bounds with no other information: $\max(\Prob(A),\Prob(B)) \le \Prob(A \cup B) \le \min(1,\Prob(A)+\Prob(B))$
        and $\max(0, \Prob(A)+\Prob(B)-1) \le \Prob(A\cap B) \le \min(\Prob(A),\Prob(B))$.
        Boole: $\Prob(\bigcup A_i) \le \sum \Prob(A_i)$.
  \item Venn bookkeeping: fill regions from the triple intersection outward (all
        three; each pair only; each single only; none); every answer is a sum of regions.
  \item Translation: ``or''/``at least one'' is union; ``and''/``both'' is
        intersection; ``not'' is complement; ``but not'' and ``only'' subtract the
        intersection; ``every policyholder has at least one'' means the union is $S$;
        ``no one has both'' means disjoint.
  \item Unions with complements: $\Prob(A\cup B) + \Prob(A \cup \cc{B}) = 1 + \Prob(A)$;
        $\Prob(A \cup \cc{B}) = 1 - \Prob(\cc{A} \cap B)$; $\Prob(\cc{A}\cup\cc{B}) = 1 - \Prob(A \cap B)$.
        Derive any of these by adding two inclusion-exclusion equations.
  \item Disjoint pair inside three events: two regions vanish; name the six that
        remain, translate each fact into a linear equation, solve outside-in.
        A fourth set disjoint from the other three adds by finite additivity:
        $\Prob(A\cup B\cup C\cup D) = \Prob(A \cup B \cup C) + \Prob(D)$.
  \item Geometric regions: $\sum_{i=1}^{N} r^i = r(1-r^N)/(1-r)$,
        $\sum_{i=1}^{\infty} r^i = r/(1-r) \le 1$. Miss probability with $N$ regions
        is $1 - r(1-r^N)/(1-r)$, floor $(1-2r)/(1-r)$; solve $r^N < c$ as
        $N > \ln c / \ln r$ and round up.
  \item Event equality: draw the outcome grid (classes $0$, $1$, $2^{+}$ per
        coordinate), mark the target, mark each description; one differing cell
        breaks equality. Rewrites that preserve equality: De Morgan,
        $F \cup G = F \cup (\cc{F}\cap G)$, absorption $F \cup (F \cap H) = F$.
  \item Never assume disjointness or independence unless stated; check that the
        regions sum to $1$ (or to the stated count).
\end{itemize}
